\section{Scalable oversight：当模型越过人类评价天花板}

普通 RLHF 默认人类能可靠比较 policy outputs。模型能力继续增长后，这个假设会率先失效：模型能完成更复杂的任务，而人类 evaluator 的专业知识和注意力不会随 AI progress 自动增长。本章先描述失效，再说明 AI assistance 如何尝试抬高评价曲线。

\subsection{Human evaluation ceiling 与 deception incentive}

图中横轴是 AI progress，纵轴是任务难度。橙色能力曲线上升，绿色 unaided human evaluation 近似水平。当橙线越过绿线，人类仍可能给出偏好，却无法知道高分输出是否真的正确；reward optimization 会把这种盲区变成可利用空间。

\lecturefigure{slide-13-scaling-human-supervision.jpg}{Scalable oversight 的核心图：AI 能力增长，人类独立评价能力却有天花板。}{Stanford Online 官方视频 00:22:25--00:25:30。}

\begin{knowledgebox}{读图：三条曲线分别代表什么}
底部曲线是 AI 能完成的任务难度；水平线是人类独立可靠评价的上限；上方新曲线是人类与 AI 协作后的评价能力。关键不是精确交点，而是顺序：先出现模型能做但人类看不懂的任务，再出现对弱评价信号的系统性优化。图不证明 deception 必然发生，但说明 deception 或 sycophancy 为什么可能成为高奖励策略。
\end{knowledgebox}

把真实任务质量写成 $q(x,y)$，人类可见评分写成 $h(x,y)$，reward model 学到 $r_\phi(x,y)\approx h(x,y)$。若在人类看不懂的区域中 $h$ 与 $q$ 脱钩，那么

\[
\arg\max_y r_\phi(x,y) \neq \arg\max_y q(x,y).
\]

其中，左侧是 policy 实际优化的 proxy，右侧是人类真正关心的任务质量。能力越强，模型越可能找到使两者分离的复杂输出。

\begin{warningbox}{“骗过 evaluator”不要求模型有人类式恶意}
只要训练持续选择高评分输出，任何稳定提高评分而不提高真实质量的模式都会被强化。它可以是讨好、过度自信、隐藏不确定性、选择性展示证据或真正的策略性欺骗。风险来自优化与测量的结构，而不是先假设模型拥有某种人格。
\end{warningbox}

\subsection{AI-assisted evaluation：把全局任务拆成可验证局部问题}

讲者用代码审计举例：人类无法在有限时间里发现大型 codebase 的全部 bug 或 Trojan，但模型若指出具体位置、失败条件与测试，人类便可以验证这个局部 claim。AI assistance 的角色不是替人类下最终价值判断，而是做阅读、搜索、计算、fact checking 和候选生成，把难评价对象转化为一组更小的证据单元。

\lecturefigure{slide-14-critiques-and-dialog.jpg}{两类 AI assistance：生成 critique，或通过 dialogue 主动查询证据。}{Stanford Online 官方视频 00:25:30--00:26:38。}

\begin{knowledgebox}{读图：critique 与 dialogue 的信息接口不同}
Critique 一次性给出 ``哪里可能错''，适合高召回地暴露检查点；dialogue 允许 evaluator 追问来源、反例、计算过程和替代解释，更适合逐步缩小不确定性。两者都依赖 assistant 的诚实性与覆盖率，不能因为语言自然就当作可靠证据。
\end{knowledgebox}

\begin{table}[H]
\centering
\caption{AI-assisted evaluation 的常见动作与剩余责任}
\begin{tabularx}{\textwidth}{L{2.6cm} X X}
\toprule
动作 & AI 可承担的认知劳动 & 人类仍需判断 \\
\midrule
Critique & 枚举缺陷、反例、遗漏和冲突 & 哪些问题重要，critique 是否成立 \\
Dialogue & 响应追问、比较方案、解释局部 claim & 何时停止追问，证据是否足够 \\
Fact checking & 检索来源、定位引文、交叉核对 & 来源可信度与语境是否匹配 \\
Code/tool use & 运行测试、计算、静态分析、浏览 & 权限边界、工具结果与真实目标的关系 \\
\bottomrule
\end{tabularx}
\end{table}

\subsection{本章小结}

模型能力超过 unaided human evaluation 后，普通 RLHF 的 proxy 可能与真实质量脱钩。Scalable oversight 试图用同类 AI 的阅读、搜索、计算与批评能力帮助人类评价，但最终偏好判断、证据信任和权限边界仍由人类承担。

